% =============================================================================
\clearpage
\appendix
\renewcommand{\thesection}{\Alph{section}}
\section{Formulario del Tema 5 (primera parte)}\label{sec:formulario}
% =============================================================================

La \cref{tab:formulario} reúne las expresiones centrales de esta parte del tema, con la referencia al enunciado en que se demuestran.

\begingroup
\centering
\captionof{table}{Formulario del Tema 5. Energía y fuerzas, primera parte ($k=1/4\pi\eo$, $\phi(\infty)=0$).}
\label{tab:formulario}
\small
\renewcommand{\arraystretch}{1.75}
\begin{tabularx}{\textwidth}{@{}>{\raggedright\arraybackslash}p{4.2cm}>{\raggedright\arraybackslash}X>{\raggedright\arraybackslash}p{2.9cm}@{}}
\toprule
Magnitud o resultado & Expresión & Referencia\\
\midrule
\multicolumn{3}{@{}l}{\textit{Energía de cargas y distribuciones}}\\
Trabajo reversible & $\Wext(A\to B)=q\big[\phi(B)-\phi(A)\big]$ & \cref{lem:trabajo}\\
Cargas puntuales (interacción) & $U_e=\dfrac12\displaystyle\sum_{i\neq j}\frac{kq_iq_j}{r_{ij}}=\frac12\sum_iq_i\phi_i$ & \cref{prop:puntuales}\\
Distribución continua (total) & $U_e=\dfrac12\displaystyle\int_V\rho\,\phi\,\dd V$,\quad $\dfrac12\displaystyle\int_S\sigma\,\phi\,\dd S$ & \cref{prop:continua}\\
\multicolumn{3}{@{}l}{\textit{Energía del campo}}\\
Energía y densidad & $U_e=\dfrac{\eo}{2}\displaystyle\int_{\mathbb R^3}E^2\,\dd V$,\quad $u_e=\dfrac{\eo}{2}E^2$,\quad $U_e\geq0$ & \cref{teo:energia-campo}, \cref{cor:positiva}\\
No aditividad & $U_e=U_1+U_2+\Uint$,\quad $\Uint=\eo\displaystyle\int\vE_1\cdot\vE_2\,\dd V=\int\rho_2\phi_1\,\dd V$ & \cref{prop:no-aditiva}\\
Corteza esférica / esfera maciza & $\dfrac{Q^2}{8\pi\eo a}$\quad / \quad$\dfrac{3Q^2}{20\pi\eo a}$ & \cref{ej:corteza}, \cref{prob:esfera-cargada}\\
\multicolumn{3}{@{}l}{\textit{Conductores}}\\
Sistema de conductores & $U_e=\dfrac12\displaystyle\sum_iQ_i\phi_i=\frac12\sum_{i,j}p_{ij}Q_iQ_j=\frac12\sum_{i,j}c_{ij}\phi_i\phi_j$ & \cref{prop:conductores}\\
Condensador & $U_e=\dfrac12QV=\dfrac12CV^2=\dfrac{Q^2}{2C}$ & \cref{cor:condensador}\\
\multicolumn{3}{@{}l}{\textit{Fuerzas}}\\
Earnshaw & no hay equilibrio estable solo con fuerzas electrostáticas & \cref{teo:earnshaw}\\
Carga constante & $\vFe=-\grad U_e|_Q$,\ \ $\tau_\theta=-\partial_\theta U_e|_Q$,\ \ $F_\xi=-\tfrac12\sum_{i,j}Q_iQ_j\,\partial_\xi p_{ij}$ & \cref{teo:fuerza-Q}\\
Potencial constante & $\vFe=+\grad U_e|_\phi$,\quad $\tau_\theta=+\partial_\theta U_e|_\phi$,\quad $\dd\Wbat=2\,\dd U_e$ & \cref{teo:fuerza-phi}\\
Equivalencia & $-\partial_\xi U_e|_Q=+\partial_\xi U_e|_\phi$ & \cref{prop:equivalencia}\\
Condensador & $F_\xi=\dfrac12V^2\dfrac{\dd C}{\dd\xi}=\dfrac{Q^2}{2C^2}\dfrac{\dd C}{\dd\xi}$ & \cref{cor:fuerza-condensador}\\
Dipolo en un campo & $\Uint=-\vp\cdot\vE$,\quad $\vF=(\vp\cdot\grad)\vE$ & \cref{prob:dipolo-fuerza}\\
\bottomrule
\end{tabularx}
\par\endgroup

% =============================================================================
\begin{thebibliography}{99}
\setlength{\itemsep}{0.25em}
\small
\bibitem{T1} D.~Martín Guerrero. \emph{Leyes básicas de la electrostática en el vacío} (Electromagnetismo I, Tema 1). Apuntes formalizados con
  seminarios prácticos. Grado en Física, Universidad de Salamanca, 2026.
\bibitem{T2} D.~Martín Guerrero. \emph{Desarrollo multipolar} (Electromagnetismo I, Tema 2). Apuntes formalizados con seminarios prácticos. Grado
  en Física, Universidad de Salamanca, 2026.
\bibitem{T3} D.~Martín Guerrero. \emph{Conductores en equilibrio electrostático} (Electromagnetismo I, Tema 3). Apuntes formalizados con seminarios
  prácticos. Grado en Física, Universidad de Salamanca, 2026.
\bibitem{T4} D.~Martín Guerrero. \emph{Materiales dieléctricos} (Electromagnetismo I, Tema 4). Apuntes formalizados con seminarios prácticos. Grado
  en Física, Universidad de Salamanca, 2026.
\bibitem{T0} D.~Martín Guerrero. \emph{Herramientas matemáticas para el electromagnetismo} (Electromagnetismo I, Tema 0). Apuntes formalizados con
  seminarios prácticos. Grado en Física, Universidad de Salamanca, 2026.
\bibitem{DFA-teoria} Departamento de Física Aplicada. \emph{Electromagnetismo I. Tema V: Energía y fuerzas}. Apuntes de teoría, §V.1--V.5. Material de
  la asignatura. Grado en Física, Universidad de Salamanca, 2025.
\bibitem{DFA-figuras} Departamento de Física Aplicada. \emph{Electromagnetismo I. Tema 5: Energía y fuerzas (figuras)}. Material de la asignatura.
  Grado en Física, Universidad de Salamanca, 2025.
\bibitem{DFA-problemas} Departamento de Física Aplicada. \emph{Electromagnetismo I. Tema 5: Energía y fuerzas (relación de problemas)}. Material de la
  asignatura. Grado en Física, Universidad de Salamanca, curso 2025-2026.
\bibitem{Guia-T5} D.~Martín Guerrero. \emph{Tema 5. Energía y fuerzas (I): formulario, métodos y fichas de problemas}. Guías de estudio de
  Electromagnetismo I. Grado en Física, Universidad de Salamanca, 2026.
\bibitem{Griffiths} D.~J. Griffiths. \emph{Introduction to Electrodynamics}. 4.ª ed. §2.4: «Work and energy in electrostatics»; §2.5.4: fuerza sobre
  conductores; §4.4.3: fuerzas sobre dieléctricos. Cambridge University Press, 2017.
\bibitem{Wangsness} R.~K. Wangsness. \emph{Electromagnetic Fields}. 2.ª ed. Cap.~7: «Electrostatic energy». John Wiley \& Sons, 1986.
\bibitem{Purcell} E.~M. Purcell y D.~J. Morin. \emph{Electricity and Magnetism}. 3.ª ed. §1.15: energía de un sistema de cargas; §3.5--3.6: energía y
  fuerzas en condensadores. Cambridge University Press, 2013.
\bibitem{Feynman} R.~P. Feynman, R.~B. Leighton y M.~Sands. \emph{The Feynman Lectures on Physics. Volume II}. Cap.~8: «Electrostatic energy». Addison-Wesley,
  1964.
\bibitem{Jackson} J.~D. Jackson. \emph{Classical Electrodynamics}. 3.ª ed. §1.11: energía potencial electrostática y densidad de energía; §1.13:
  coeficientes de capacidad. John Wiley \& Sons, 1999.
\bibitem{Kellogg} O.~D. Kellogg. \emph{Foundations of Potential Theory}. Cap.~VIII: teoremas del valor medio y principio del máximo. Berlín: Springer, 1929.
\bibitem{Earnshaw} S.~Earnshaw. «On the nature of the molecular forces which regulate the constitution of the luminiferous ether». \emph{Transactions of the
  Cambridge Philosophical Society} 7 (1842), 97--112.
\end{thebibliography}
